\section{Глава~\ref{sec:pain}. Боль}

Датчики боли, два провода, ворота в спинном мозге, нисходящая ручка громкости, сенсибилизация и захват внимания измерены прямо; формулы ворот и сенсибилизации --- упрощённые модели книги; правило, по которому машина выбирает громкость тревоги, --- гипотеза.

{\footnotesize
\setlength{\tabcolsep}{3pt}\renewcommand{\arraystretch}{1.15}
\begin{longtable}{>{\raggedright}p{0.5cm}>{\raggedright}p{4.0cm}>{\raggedright}p{5.6cm}>{\raggedright}p{2.3cm}>{\raggedright\arraybackslash}p{1.7cm}}
\toprule
№ & Утверждение & Опыт и что измерено & Источник & Статус \\
\midrule\endfirsthead
\toprule
№ & Утверждение & Опыт и что измерено & Источник & Статус \\
\midrule\endhead
1 & Высокий порог: датчики боли срабатывают при нагреве выше примерно $43^\circ$ & Запись одиночных волокон. Кожа обезьяны: медианный порог медленных (C) датчиков $41^\circ$, быстрых датчиков второго типа $46^\circ$, первого типа --- выше $53^\circ$. Кожа человека: C-датчики, отвечающие и на давление, $40.7^\circ$, не отвечающие на давление --- $48^\circ$. Число $43^\circ$ --- середина этого диапазона, а не единый порог. & \cite{Treede1995heat, Weidner1999cfib} & измерено \\
2 & После повреждения датчики не глохнут, а становятся чувствительнее & Лёгкий ожог кожи ($50^\circ$, $100\,\text{с}$): через $5$--$10$~мин порог боли у людей и порог C-датчиков у обезьяны снижены на $2$--$6^\circ$; у других кожных датчиков такого сдвига нет. Слабое привыкание прямо не сравнивалось; у быстрых датчиков второго типа ответ после сильного нагрева, наоборот, подавлен. & \cite{LaMotte1982hyper, Treede1995heat} & измерено \\
3 & Два провода: быстрый $5$--$30$~м/с, медленный $0.5$--$2$~м/с; время прихода $t=L/v$~\eqref{eq:paintime} & Скорость проведения: быстрые датчики боли обезьяны в среднем $15$ и $25$~м/с (два типа); C-датчики человека $0.8$--$1.0$~м/с. При $L=1$~м это десятки миллисекунд и около секунды. & \cite{Treede1995heat, Weidner1999cfib} & измерено \\
4 & Боль приходит дважды: сначала быстрая и точная, потом тупая и долгая & Время реакции на нагрев предплечья человека: от $1100$ до $400\ms$ с ростом температуры; сильный нагрев волосистой кожи даёт двойную боль, ладони --- нет. Первую боль могут нести только волокна быстрее $6$~м/с --- по латентности ей соответствуют быстрые датчики второго типа, которых в ладони нет. Раздел ролей («где» и «насколько плохо») --- толкование. & \cite{CampbellLaMotte1983first, Treede1995heat} & измерено \\
5 & Ворота: выходной нейрон спинного мозга получает боль и касание, тормозящий посредник касанием возбуждается (рис.~\ref{fig:gate}) & Схема ворот предложена как теория. Мыши, генетическая метка и удаление групп нейронов: возбуждающие нейроны с соматостатином нужны для механической боли; тормозящие нейроны с динорфином не дают входу от датчиков касания дойти до них --- без этих нейронов лёгкое касание вызывает боль. & \cite{MelzackWall1965, Duan2014} & измерено \\
6 & Касание приглушает боль & Человек, лазерные импульсы боли на руке: одновременное касание снижает оценку боли и способность различать энергию импульсов; эффект слабеет с расстоянием между точкой касания и точкой боли внутри одного кожного сегмента. & \cite{Mancini2014touch} & измерено \\
7 & Сильная боль сама выключает тормозящего посредника, и ворота распахиваются & Это часть исходной схемы ворот. В приведённых опытах прямо не показано: работа на мышах описывает, как ворота не пускают касание в канал боли, а не выключение посредника болью. & \cite{MelzackWall1965} & гипотеза \\
8 & Ворота~\eqref{eq:gate}: порог $0.18$ без касания, $0.86$ с касанием, $0.11$ при $g=2$; при $\nu=1$ касание снижает выход с $1$ до $0.25$, при $\nu=2$ не действует; касание само не проходит (рис.~\ref{fig:gatecurves}) & Расчёт по \texttt{models/\allowbreak pain\_\allowbreak gate.py} при весах из текста воспроизводит все эти числа. & \texttt{models/\allowbreak pain\_\allowbreak gate.py} & модель \\
9 & Нисходящие пути из ствола меняют усиление ворот в обе стороны & Крыса, продолговатый мозг, запись при нагреве хвоста: одни нейроны разряжаются перед отдёргиванием («вкл»), другие замолкают («выкл»); слабая стимуляция этой области подавляет рефлекс. Обзор: те же цепи и облегчают, и тормозят передачу боли, опиоиды действуют через них. & \cite{Fields1983rvm, Fields2004} & измерено \\
10 & Ожидание облегчения приглушает боль через опиоидный канал & Люди с острой болью: у тех, кому ожидание облегчения снизило боль, блокада опиоидных рецепторов вернула боль к уровню остальных; у остальных блокада боль не меняла. & \cite{Levine1978} & измерено \\
11 & Громкость тревоги мозг выбирает по выгоде прерывания & Опыта, в котором измерено правило выбора громкости, нет. & --- & гипотеза \\
12 & Сенсибилизация: после повреждения выход ворот растёт, порог падает, отчасти за счёт спинного мозга & Крыса: после повреждения кожи порог сгибательного рефлекса падает, а ответ растёт; электрофизиологический анализ показывает, что часть роста возникает в самом спинном мозге. Время спада «от часов до дней» в этой работе не измерялось. & \cite{Woolf1983} & измерено \\
13 & Сенсибилизация~\eqref{eq:sensitization} --- положительная обратная связь; при сильной петле тревога может защёлкнуться после заживления & Следует из модели~\eqref{eq:sensitization} (задача в конце главы). Опыта, показывающего, что стойкая боль после заживления --- защёлкнувшаяся петля именно такого вида, нет. & вывод в главе & гипотеза \\
14 & Боль --- отрицательная награда, и обучение на ней следует ошибке временных разностей & Томограф, люди, цепочки предвестников боли: активность вентрального стриатума и передней островковой коры совпадает с сигналами обучения по временным разностям. Обезьяна: часть \nt{ДОФА}-нейронов возбуждается от неприятного обдува и его предвестника, другая часть тормозится. & \cite{Seymour2004, Matsumoto2009} & измерено \\
15 & Боль прерывает текущую работу & Люди, электрический болевой стимул и звуковые пробы: ответ на пробу через $100\ms$ после начала боли заметно замедлен, через $1.5\,\text{с}$ разницы с контролем уже нет. Обзор сводит такие опыты в модель захвата внимания. & \cite{Crombez1996disrupt, EcclestonCrombez1999} & измерено \\
16 & Отдёргивание замыкается в спинном мозге & Крыса: нейроны глубоких слоёв заднего рога повторяют пространственный узор рефлекса отдёргивания отдельных мышц; их не удаётся возбудить антидромно из шейного отдела, то есть это спинальные посредники. & \cite{Schouenborg1995wr} & измерено \\
\bottomrule
\end{longtable}}
